
% A propagação de ondas de radio são rigorosamente reguladas por orgãos
% competentes à homologação de dispositivos de comunicação, onde há regras de funcionamento
% definidas por cada país. No caso do Brasil
% a ANATEL (Agência Nacional de Telecomunicações), é a responsável 
% por regular, fiscalizar e homologar dispositivos com o foco de estabelecer redes de comunicação.
% Na Europa, cada país há um orgão interno voltado à regulação de aparelhos de comunicação, 
% mas por convenção, todos os países-membros da união europeia desenvolvem suas regulações 
% junto da BEREC (Body of European Regulators for Electronic Communications).

% O \textit{duty cycle} é a relação entre o tempo em que um dispositivo 
% ocupa um canal de rádio e o tempo total de observação. No plano europeu 
% EU868 (863--870~MHz), há a imposição de \textit{duty cycle} máximo 
% de 1\%, o que equivale a um orçamento de transmissão de apenas 
% 36~segundos por hora, ou 864~segundos por dia \cite{jebril2018overcoming}.

% O \textit{dwell time} é o tempo máximo de ocupação contínua de uma 
% frequência por transmissão. No Brasil, o padrão regulamentado pela 
% ANATEL para redes \gls{lorawan} é o AU915, operando na faixa de 
% 902--928~MHz. Diferentemente do EU868, o AU915 não impõe restrição 
% de \textit{duty cycle}; em seu lugar, o Ato nº~14.448, de 4 de dezembro 
% de 2017~\cite{anatel2017} estabelece que, para sistemas de salto em 
% frequência operando nas faixas de 902--907,5~MHz e 915--928~MHz com 
% largura de faixa de canal inferior a 250~kHz, caso do \gls{lora} com 
% 125~kHz, o tempo médio de ocupação de qualquer radiofrequência não 
% deve ser superior a \textbf{0,4~segundos} em um intervalo de 
% 14~segundos, com uso mínimo de 35~canais de salto. O AU915 utiliza 
% 64~canais de \textit{uplink}, satisfazendo amplamente essa condição.

% Valores mais altos de \gls{sf} aumentam a sensibilidade do receptor e o 
% alcance da transmissão, porém elevam proporcionalmente o tempo que o 
% sinal permanece no ar (\textit{time on air}), impactando diretamente 
% o consumo energético e a capacidade da rede~\cite{adam2022}.

% O tempo no ar por pacote ($T_{\text{pacote}}$) é calculado a partir 
% do \gls{sf}, da \gls{bw}, do \gls{cr} e do tamanho do \textit{\gls{payload}}, 
% conforme a formulação oficial do \gls{lora}. 
% Como explicado anteriormente, o \gls{cr} expressa a proporção de bits 
% redundantes adicionados para tolerância a erros; o valor 4/5 representa 
% o menor \textit{overhead} de redundância disponível no padrão 
% \gls{lora}~\cite{adam2022}.

Ao redor do mundo, protocolos de transmissão que utilizam ondas de rádio 
como meio de propagação de sinal, tem seu funcionamento regulado por 
órgãos competentes à homologação de dispositivos de comunicação.

No caso do Brasil, a ANATEL (Agência Nacional de Telecomunicações) é a 
responsável por regular, fiscalizar e homologar dispositivos com o foco de 
estabelecer redes de comunicação. 
Na Europa, cada país tem um órgão interno voltado à regulação de aparelhos 
de comunicação, mas por convenção, todos os países-membros da União Europeia 
desenvolvem suas regulações junto da BEREC (Body of European 
Regulators for Electronic Communications). 

O AU915 é a restrição imposta pela ANATEL no ato nº 14.448, estabelecida 
em 4 de dezembro de 2017~\cite{anatel2017}, em que se adota o 
padrão australiano de radiocomunicação. 
Por mais que o Brasil adotou o padrão americano de distribuição de canais \gls{ism}, 
o AU915 foi escolhido por conta da já existente implantação de infraestrutura 
de rádio que ocupavam parcialmente as faixas sub 915Mhz, portanto o padrão 
americano (US915) que funciona em uma faixa de frequência entre 902 e 928 Mhz, 
tornou-se inadequado para implementação em território brasileiro.

O \gls{ism}, foi desenvolvido em Atlanta em 1947, 
elaborado pela FCC (Federal Communications Comission), 
com o objetivo de delimitar espectros de radiofrequência 
para as mais diversas finalidades. 
Em 1985, o FCC autorizou o uso de livre de espectros 
de rádio em bandas ISM sem a necessidade de licenciamento para 
utilização em determinados espectros de frequência.

Entre as regulações implementadas pelos mais diversos países ao redor do mundo, dois termos ganham grande relevância. 
O \textit{duty cycle} é a relação entre o tempo em que um dispositivo ocupa um canal de rádio e o tempo total de observação. 
E o \textit{dwell time} é o tempo máximo de ocupação contínua de uma frequência por transmissão.

No plano europeu EU868 (863–870 MHz), por exemplo, há a imposição de duty cycle máximo de 1\%, o que equivale a um orçamento de transmissão de apenas 36 segundos por hora, ou 864 segundos por dia (JEBRIL et al., 2018). 

Tendo o AU915 operando entre as faixas de frequência de 915 e 928 MHz, a norma brasileira não impõe restrição de duty cycle, em seu lugar, o Ato nº 14.448 estabelece que, para sistemas de salto em frequência operando nas faixas de 902–907,5 MHz e 915–928 MHz com largura de faixa de canal inferior a 250 kHz, caso do LoRa com 125 kHz, o tempo médio de ocupação de qualquer radiofrequência não deve ser superior a 0,4 segundos em um intervalo de 14 segundos, com uso mínimo de 35 canais de salto. 

\begin{table}[H]
  \centering
  \caption{Uplink: Canais de Frequência}
  \label{tab:uplink_freq_channel}
  \begin{tabular}{|c|c|}
    \hline
    \textbf{Canal} & \textbf{Frequência Inicial (Mhz)} \\ \hline
    0  & 915,2 \\ \hline
    1  & 915,4  \\ \hline
    2  & 915,6  \\ \hline
    … &  …  \\ \hline
    67 & 927,8 \\ \hline
  \end{tabular}
\end{table}

\begin{table}[H]
  \centering
  \caption{Downlink: Canais de Frequência}
  \label{tab:downlink_freq_channel}
  \begin{tabular}{|c|c|}
    \hline
   \textbf{Canal} & \textbf{Frequência Inicial (Mhz)} \\ \hline
    0  & 923,3 \\ \hline
    1  & 923,8  \\ \hline
    2  & 924,3  \\ \hline
    … &  …  \\ \hline
    7 & 927,5 \\ \hline
  \end{tabular}
\end{table}

O AU915 utiliza 64 canais de uplink e 8 canais de downlink em que tais canais se distribuem de entre os espectros de frequência delimitados pela restrição regulatória. No caso de uplink, como mostrado pela tabela~\ref{tab:uplink_freq_channel}, os canais são numerados entre 0 e 63 e distribuídos com espaçamentos de 200 Mhz, iniciando em 915,2 Mhz até 927,8 Mhz. E no caso de Downlink, como mostrado na tabela~\ref{tab:downlink_freq_channel}, os canais são numerados de 0 a 7, e distribuídos entre 923,3 e 927,5 Mhz com espaçamentos de 500 Mhz. \cite{lora_alliance_rp002_105}.

